\chapter*{Nota sobre las fuentes y las cifras}
\addcontentsline{toc}{chapter}{Nota sobre las fuentes y las cifras}
\markboth{Nota sobre las fuentes}{Nota sobre las fuentes}

Las estadísticas cubanas son un problema de tres capas, y un libro que las cite
sin decirlo engaña a su lector por aritmética en vez de por argumento.

La primera capa es corriente. La Oficina Nacional de Estadística e Información
(ONEI) publica un \emph{Anuario Estadístico} detallado, internamente
consistente y, para los estándares de la región, razonablemente competente. Las
series de salud y educación en particular las recoge un aparato administrativo
que llega a cada municipio, y la cobertura del registro civil es casi completa,
lo que es más de lo que puede decirse de la mayoría de los países del nivel de
ingreso de Cuba. Donde este libro cita esperanza de vida, mortalidad infantil,
matrícula escolar o número de médicos, la fuente suele ser la ONEI y la cifra
suele ser fiable en el sentido estricto de que el Estado midió lo que dice haber
medido.

La segunda capa es el tipo de cambio, y corrompe todo lo monetario. Durante casi
todo el período que aquí se cubre, el Estado cubano mantuvo una tasa oficial de
un peso por dólar para la contabilidad empresarial mientras los hogares cambiaban
a veinticuatro o veinticinco, y después de 2021 a tasas que en el mercado
informal superaron las trescientas. Cualquier cifra denominada en pesos cubanos y
convertida a la tasa oficial ---el PIB, el PIB por habitante, los salarios, el
valor del salario social, el costo de un subsidio--- no es entonces una medición
sino una convención contable, y las comparaciones entre países construidas sobre
ella carecen casi por completo de sentido. El PIB por habitante declarado de
Cuba la ha situado durante décadas en la banda media-alta de América Latina
mientras su canasta de consumo parecía la de la banda baja. Ambos hechos son
reales; la contradicción está en la conversión. Este libro da las cifras en pesos
con la tasa adjunta, prefiere las magnitudes físicas a las monetarias siempre que
exista una magnitud física ---toneladas de azúcar, barriles de petróleo,
megavatios, médicos por mil habitantes, calorías diarias--- y trata cualquier
cifra única del PIB cubano como un argumento antes que como un dato.

La tercera capa es política, y corre en las dos direcciones. El Estado no
publica, o publica tarde y de forma parcial, las series que lo incomodan: la
magnitud de la contracción de 1993--94 se reconstruyó años después, la población
penal no se divulga, el conteo de detenciones políticas lo producen observadores
de la oposición y nadie más, y después de 2020 varias series rutinarias del
\emph{Anuario} sencillamente dejaron de aparecer. Frente a esto, la literatura
del exilio tiene sus propias inflaciones, algunas muy grandes y muy duraderas
---los veinte mil muertos atribuidos a los últimos años de Batista, cifra que
nació en una revista y que jamás ha sido respaldada por nada parecido a un
conteo, es el ejemplo más conocido, y todavía hoy la repiten personas que jamás
aceptarían una cifra del gobierno cubano con la misma evidencia detrás---. Donde
en este libro aparece una cifra en disputa, aparece como un rango con el
desacuerdo nombrado, y donde la respuesta honesta es que nadie lo sabe, el texto
dice que nadie lo sabe.

De ahí se siguen tres convenciones prácticas.

Las cifras afirmadas a partir de una fuente secundaria que no se ha rastreado
hasta una primaria van marcadas en el manuscrito y listadas en
\texttt{verify.md}, en el directorio fuente del libro. La marca no imprime nada,
de modo que no ensucia la página; existe para que la lista de afirmaciones sin
verificar sea extraíble mecánicamente y no cuestión de memoria.

Donde se da un rango, sus extremos son las dos estimaciones publicadas más
defendibles y no los extremos de lo que se ha llegado a afirmar. Una cifra
sostenida únicamente por una parte interesada en ella, y nunca de forma
independiente, se describe así en vez de promediarse dentro de un rango como si
fuera evidencia.

Donde el libro enuncia una afirmación causal ---que el embargo le costó a Cuba
tanto, que la reforma monetaria de 2021 causó tal inflación--- enuncia qué
tendría que ser cierto para que se sostenga, porque en esta materia casi toda
afirmación causal la disputa alguien cuya objeción no es frívola. Las dos
preguntas de fondo, cuánto de la pobreza cubana es el embargo y cuánto es el
modelo, y cuánto del logro social cubano es el socialismo y cuánto es lo que
cualquier Estado decidido puede comprar con un subsidio externo suficientemente
grande, no quedan resueltas en este libro. Se mantienen, eso sí, a la vista,
porque el diseño de la Parte III tiene que funcionar bajo cualquiera de las dos
respuestas.
